\section{The \texttt{Shell\_h} McXtrace Component}
Single Pore as part of the Silicon Pore Optics (SPO) as envisioned for the ATHENA+ space telescope.

\subsection*{Identification}
\begin{itemize}
  \item \textbf{Author:} Erik B Knudsen and Desiree D. M. Ferreira
  \item \textbf{Origin:} DTU Physics, DTU Space
  \item \textbf{Date:} Feb. 2016
\end{itemize}

\subsection*{Description}
A single shell is simulated. The top and bottom are curved cylindrically azimuthally, and according to the Wolter I optic lengthwise (sagitally). This is the hyperbolic part. The azimuthal curvature is defined by the radius parameters.

To intersect the Wolter I plates we take advantage of the azimuthal symmetry and only consider the radial component of the photon's wavevector.

Imperfect mirrors may be modelled using one of 4 models. In all cases the surface normal of the mirror at the ideal mirror intersection point is perturbed before the exit vector is computed. 1. Longitudinal 1D. A perturbation angle is chosen from a uniform distribution with width waviness. 2. Isotropic 2D. The surface normal is perturbed by choosing an angle on a disc with radius waviness 3. Externally measured/computed data. We interpolate in a data-file consisting of blocks of dtheta/theta with 1 block per energy. dtheta is a sampled angle offset from the nominal Fresnel grazing angle theta. 4. Double gaussian. dtheta is chosen from one of two gaussian distributions. Either specular or off-specular, where the widths (sigmas) are given by the tables in the file "wave\_file". If the off-specular case the behaviour is similar to 2D uniform case.

In the case of 3, the format of the data file should be: \#e\_min=0.1 \#e\_max=15 \#e\_step=0.01 \#theta\_min=0.01 \#theta\_max=1.5 \#theta\_step=0.01 \#dtheta\_min=-0.02 \#dtheta\_max=0.02 \#dtheta\_step=0.001

\begin{verbatim}
1.0  0.9  0.8  0.75  ...
0.99 0.89 0.79 0.749 ...
\end{verbatim}

... \#block 2 (energy data point 2)

\begin{verbatim}
1.0  0.9  0.8  0.75  ...
0.99 0.89 0.79 0.749 ...
\end{verbatim}

...

I.e. one 2D data block per energy data point where rows represent the steps in nominal incident angle, and columns represent the sampled granularity of the off-specular scattering.

Example: Shell\_h( radius\_m=0.535532, radius\_h=0.533113, zdepth=0.5, Z0=FL, yheight=1e-2, R\_d=1)

\subsection*{Input parameters}
Parameters in \textbf{boldface} are required; the others are optional.

\begin{longtable}{p{0.22\textwidth}p{0.12\textwidth}p{0.46\textwidth}p{0.14\textwidth}}
\toprule
\textbf{Name} & \textbf{Unit} & \textbf{Description} & \textbf{Default} \\
\midrule
\endhead
\textbf{radius\_m} & m & Ring radius of the upper (reflecting) plate of the pore at the intersection with the parabolic section. &  \\
\textbf{radius\_h} & m & Ring radius of the upper (reflecting) plate of the pore at the edge closest to the focal point. &  \\
\textbf{Z0} & m & distance between intersection plane and the focal spot( essentially the focal length). &  \\
\textbf{yheight} & m & Height of the pore. (Thus the inner radius is radius\_\{m,h\}-yehight &  \\
chamferwidth & m & Width of side walls. & 0 \\
gap & m & gap between intersection with parabolic section and actual plate. & 0 \\
zdepth &  &  & 0 \\
mirror\_reflec &   & Data file containing reflectivities of the reflector surface (TOP). & "" \\
bottom\_reflec &   & Data file containing reflectivities of the bottom surface (BOTTOM). & "" \\
wave\_file &  &  & "" \\
R\_d &   & Default reflectivity value to use if no reflectivity file is given. Useful f.i. is one surface is reflecting and the others absorbing. & 1 \\
wave\_model &   & Flag to choose waviness model. 1. longitudinal uniform, 2. 2D-uniform, 3. lorentzian sagittal, 4. double gaussian sagittal. See above for details. & 0 \\
waviness & rad & Waviness of the pore reflecting surface. The slope error is assumed to be uniformly distributed in the interval [-waviness:waviness]. & 0 \\
verbose &   & If !=0 output extra info during simulation. & 0 \\
\bottomrule
\end{longtable}

\subsection*{Links}
\begin{itemize}
  \item Component source code found in file \texttt{Shell\_h.comp}.
\end{itemize}
\IfFileExists{astrox/Shell_h_static.tex}{\input{astrox/Shell_h_static.tex}}{}